% ============================================================
%  REPA – Appendix / Backup Slides (English) with Vietnamese notes
% ============================================================

% ------------------------------------------------------------------
% APPENDIX SLIDE 1 – ROADMAP
% ------------------------------------------------------------------
\begin{frame}[t]
  \framesubtitle{Appendix}
  \frametitle{Backup slide roadmap}
  \begin{columns}[T,totalwidth=\textwidth]
    \column{0.48\textwidth}
    \begin{enumerate}
      \item DDPM vs stochastic interpolants
      \item Velocity and score relations
      \item CKNNA and linear probing
      \item DiT / SiT block details
      \item DiT analysis (not only SiT)
    \end{enumerate}
  \column{0.48\textwidth}
    \begin{enumerate}\setcounter{enumi}{5}
      \item More ablations (\(\lambda\), timesteps)
      \item Hyperparameters and compute
      \item Detailed quantitative tables
      \item ImageNet 512\(\times\)512 extension
      \item Text-to-image extension
      \item Feature maps and extra samples
    \end{enumerate}
  \end{columns}
  \note{
    Đây là danh sách các backup slide. Các bạn có thể hỏi về bất kỳ
    chủ đề nào, mình sẽ chuyển đến slide tương ứng.
  }
\end{frame}

% ------------------------------------------------------------------
% APPENDIX SLIDE 2 – DDPM vs STOCHASTIC INTERPOLANTS
% ------------------------------------------------------------------
\begin{frame}[t]
  \framesubtitle{Appendix}
  \frametitle{DDPM vs stochastic interpolants}
  \begin{columns}[T,totalwidth=\textwidth]
    \column{0.48\textwidth}
    \textbf{DDPM view}
    \begin{itemize}
      \item Discrete-time forward corruption with schedule \(\beta_t\)
      \item Learn reverse transitions \(p(x_{t-1} \mid x_t)\)
      \item Common objective predicts noise \(\epsilon\)
    \end{itemize}
    \begin{align*}
      \mathcal{L}_{\text{simple}}
      = \mathbb{E}_{x^\ast,\epsilon,t}\bigl\|
        \epsilon - \epsilon_\theta(x_t, t)
      \bigr\|_2^2
    \end{align*}
  \column{0.48\textwidth}
    \textbf{Stochastic interpolant view}
    \begin{itemize}
      \item Continuous path between data and Gaussian noise
      \item Learn velocity or score over continuous \(t\)
      \item Cleaner interface for flow-matching models (SiT)
    \end{itemize}
    \begin{align*}
      x_t = \alpha_t x^\ast + \sigma_t\,\epsilon,\quad
      \dot{x}_t = v(x_t, t)
    \end{align*}
    \vspace{0.12cm}
    {\small REPA is compatible with both views because it attaches to
    hidden representations, not to a specific sampler or parameterization.}
  \end{columns}
  \note{
    So sánh hai framework. DDPM: thời gian rời rạc, học phân phối
    ngược $p(x_{t-1} | x_t)$, loss predict noise $\epsilon$.

    Stochastic interpolants: thời gian liên tục, nội suy trực tiếp
    giữa data và noise, học velocity field. SiT dùng framework này.

    Điểm quan trọng: REPA hoạt động với cả hai vì nó chỉ ``nhìn''
    vào hidden state, không quan tâm model predict noise hay velocity.
    Paper thực nghiệm trên cả DiT (DDPM-style) và SiT (interpolant-style).
  }
\end{frame}

% ------------------------------------------------------------------
% APPENDIX SLIDE 3 – VELOCITY & SCORE RELATIONS
% ------------------------------------------------------------------
\begin{frame}[t]
  \framesubtitle{Appendix}
  \frametitle{Velocity and score relations}
  \textbf{Velocity field:}
  \begin{align*}
    v(x, t) &= \dot{\alpha}_t\,\mathbb{E}[x^\ast \mid x_t{=}x]
               + \dot{\sigma}_t\,\mathbb{E}[\epsilon \mid x_t{=}x]
  \end{align*}
  \textbf{Velocity loss:}
  \begin{align*}
    \mathcal{L}_{\text{velocity}}(\theta)
    &= \mathbb{E}_{x^\ast,\epsilon,t}\bigl\|
      v_\theta(x_t, t) - \dot{\alpha}_t x^\ast - \dot{\sigma}_t\,\epsilon
    \bigr\|_2^2
  \end{align*}
  \textbf{Score function:}
  \begin{align*}
    s(x, t) &= -\sigma_t^{-1}\,\mathbb{E}[\epsilon \mid x_t{=}x]
  \end{align*}
  \textbf{Score from velocity conversion:}
  \begin{align*}
    s(x, t) = \sigma_t^{-1} \cdot
    \frac{\alpha_t\,v(x, t) - \dot{\alpha}_t\,x}
         {\dot{\alpha}_t\,\sigma_t - \alpha_t\,\dot{\sigma}_t}
  \end{align*}
  \vspace{0.1cm}
  \callout{\textbf{Why these equations matter for REPA.} The denoising math stays intact.
  REPA only adds feature-level supervision on the same hidden states used to compute these objectives.}
  \note{
    Chi tiết công thức toán. Velocity field $v(x,t)$ là tổ hợp tuyến
    tính của conditional expectation của $x^*$ và $\epsilon$.

    Loss velocity: minimize khoảng cách L2 giữa velocity dự đoán
    và ground truth velocity.

    Score function $s(x,t)$: liên quan đến gradient của log density.
    Có thể convert từ velocity sang score qua công thức chuyển đổi.

    Điểm mấu chốt: REPA không thay đổi bất kỳ công thức nào ở đây.
    Nó chỉ thêm supervision ở mức feature lên cùng hidden state
    mà model dùng để tính velocity/score.
  }
\end{frame}

% ------------------------------------------------------------------
% APPENDIX SLIDE 4 – CKNNA & LINEAR PROBING
% ------------------------------------------------------------------
\begin{frame}[t]
  \framesubtitle{Appendix}
  \frametitle{What are CKNNA and linear probing measuring?}
  \begin{columns}[T,totalwidth=\textwidth]
    \column{0.52\textwidth}
    \textbf{Linear probing}
    \begin{itemize}
      \item Freeze the representation.
      \item Train a linear classifier on ImageNet labels.
      \item Higher accuracy \(\Rightarrow\) stronger semantic separability.
    \end{itemize}
    \vspace{0.16cm}
    \textbf{CKNNA (Centered Kernel Nearest-Neighbor Alignment)}
    \begin{itemize}
      \item A relaxed alignment metric related to CKA.
      \item Focuses on mutual nearest neighbors rather than
            all pairwise relationships.
      \item More forgiving when comparing representations from
            different model families.
    \end{itemize}
  \column{0.44\textwidth}
    \textbf{CKA formula:}
    \begin{align*}
      \mathrm{CKA}(K, L)
      = \frac{\mathrm{HSIC}(K, L)}
             {\sqrt{\mathrm{HSIC}(K, K)\,\mathrm{HSIC}(L, L)}}
    \end{align*}
    \begin{beamercolorbox}[wd=\linewidth,rounded=true,sep=0.18cm]{block body}
      \textbf{How the paper uses them}
      \begin{itemize}
        \item Linear probing: measures the \textbf{semantic gap}
              between diffusion and SSL features
        \item CKNNA: measures \textbf{representational alignment}
              between two feature spaces
      \end{itemize}
    \end{beamercolorbox}
    \vspace{0.14cm}
    {\small Both metrics improve significantly with REPA,
    confirming that alignment is happening.}
  \end{columns}
  \note{
    Giải thích hai metric dùng trong paper.

    Linear probing: đóng băng representation, train linear classifier
    trên ImageNet. Accuracy cao = features chứa nhiều thông tin ngữ
    nghĩa. Paper dùng để đo ``semantic gap''.

    CKNNA: metric đo sự tương đồng giữa hai representation space.
    Dựa trên CKA (Centered Kernel Alignment) nhưng relax hơn —
    chỉ xét mutual nearest neighbors thay vì tất cả cặp.

    CKA dùng HSIC (Hilbert-Schmidt Independence Criterion) để đo
    sự phụ thuộc giữa hai kernel matrix K và L.

    Paper dùng linear probing để đo semantic quality, CKNNA để đo
    alignment với DINOv2. Cả hai đều cải thiện đáng kể với REPA.
  }
\end{frame}

% ------------------------------------------------------------------
% APPENDIX SLIDE 5 – DiT/SiT BLOCK DETAILS
% ------------------------------------------------------------------
\begin{frame}[t]
  \framesubtitle{Appendix}
  \frametitle{DiT / SiT block details}
  \begin{columns}[T,totalwidth=\textwidth]
    \column{0.42\textwidth}
    \FigNine[width=\linewidth]
  \column{0.54\textwidth}
    \begin{itemize}
      \item Patchified latent tokens go through repeated
            \textbf{attention + MLP blocks}.
      \item \textbf{AdaLN-zero} style modulation injects timestep
            and class condition information.
      \item REPA taps one \textbf{intermediate hidden state} and sends it
            to a separate projection MLP.
      \item The projection head is a \textbf{training-time adapter}
            into the external encoder feature space.
      \item At inference: projector discarded, model runs at
            \textbf{original speed}.
    \end{itemize}
    \vspace{0.12cm}
    \callout{\textbf{Model sizes.}\\
    DiT/SiT-B: 12 blocks, 768 hidden dim\\
    DiT/SiT-L: 24 blocks, 1024 hidden dim\\
    DiT/SiT-XL: 28 blocks, 1152 hidden dim}
  \end{columns}
  \note{
    Chi tiết kiến trúc DiT và SiT.

    Input: ảnh qua VAE encoder thành latent 32x32x4, rồi patchify
    thành sequence of tokens (ví dụ 256 token cho patch size 2).

    Mỗi block gồm: multi-head self-attention, rồi MLP. Cả hai đều
    dùng AdaLN-zero — adaptive layer norm modulated bởi timestep $t$
    và class label $c$.

    REPA ``tap'' vào output của một block sớm (thường block 6-8
    cho XL với 28 blocks). Output này đi qua projection MLP 3 layer
    để ánh xạ vào feature space của DINOv2.

    Khi inference: bỏ projector, model chạy giống hệt gốc — không
    có overhead nào. Đây là lợi thế lớn so với các phương pháp cần
    thay đổi architecture vĩnh viễn.

    Các model size: B (12 block), L (24 block), XL (28 block).
  }
\end{frame}

% ------------------------------------------------------------------
% APPENDIX SLIDE 6 – DiT analysis
% ------------------------------------------------------------------
\begin{frame}[t]
  \framesubtitle{Appendix}
  \frametitle{The same behavior also appears in DiT}
  \begin{columns}[T,totalwidth=\textwidth]
    \column{0.47\textwidth}
    \FigTen[width=\linewidth]
  \column{0.49\textwidth}
    \begin{itemize}
      \item DiT hidden states also show a \textbf{semantic gap}
            relative to DINOv2.
      \item Alignment is present but \textbf{weak}, just like in SiT.
      \item The paper's argument generalizes across denoising
            transformer families — it is \textbf{not a SiT-only artifact}.
    \end{itemize}
    \vspace{0.14cm}
    \callout{REPA applied to DiT-XL/2: matches vanilla DiT's 7M-step
    FID (9.6) in only 850K iterations — \(\mathbf{8\times}\) faster.}
    \vspace{0.14cm}
    {\small This shows the representation bottleneck is a general
    property of denoising transformers, not specific to the SiT
    velocity parameterization.}
  \end{columns}
  \note{
    Paper không chỉ thử trên SiT mà cả DiT. Figure này cho thấy
    DiT cũng có hiện tượng tương tự: hidden state có semantic
    information nhưng alignment với DINOv2 yếu.

    Khi áp dụng REPA cho DiT-XL/2: chỉ cần 850K iteration để đạt
    FID 9.6 — bằng với vanilla DiT sau 7 triệu iteration.
    Tức nhanh hơn 8 lần.

    Kết luận: representation bottleneck là tính chất chung của
    denoising transformer, không phụ thuộc vào cách parameterize
    (noise prediction hay velocity prediction).
  }
\end{frame}

% ------------------------------------------------------------------
% APPENDIX SLIDE 7 – MORE ABLATIONS
% ------------------------------------------------------------------
\begin{frame}[t]
  \framesubtitle{Appendix}
  \frametitle{More ablations: timesteps, target encoders, and \(\lambda\)}
  \begin{columns}[T,totalwidth=\textwidth]
    \column{0.58\textwidth}
    \papercrop[width=\linewidth]{10}{70bp 320bp 60bp 88bp}
    \vspace{0.08cm}
    {\small The representation gap (measured by CKNNA) improves across
    \emph{all} noise levels when REPA is applied.}
  \column{0.38\textwidth}
    \begin{itemize}
      \item REPA improves the representation gap across multiple
            noise levels, not just one regime.
      \item Gains are not specific to DINOv2; MAE and MoCov3
            also show the trend.
    \end{itemize}
    \vspace{0.14cm}
    \textbf{\(\lambda\) sensitivity:}
    \vspace{0.08cm}

    \begin{tabular}{C{1.4cm}C{1.0cm}C{1.0cm}C{1.0cm}}
      \toprule
      \(\lambda\) & 0.25 & 0.5 & 1.0 \\
      \midrule
      FID \metricdown & 8.6 & 7.9 & 7.8 \\
      \bottomrule
    \end{tabular}
    \vspace{0.10cm}

    {\small Performance is fairly robust once \(\lambda \geq 0.5\).
    The default \(\lambda = 0.5\) is a safe choice.}
  \end{columns}
  \note{
    Ablation chi tiết. Hình bên trái cho thấy CKNNA alignment
    cải thiện ở mọi mức noise level khi dùng REPA — không chỉ
    ở noise thấp hay noise cao.

    Bên phải: sensitivity với $\lambda$. FID khá stable:
    $\lambda = 0.25$: FID 8.6, $\lambda = 0.5$: 7.9, $\lambda = 1.0$: 7.8.
    Tức là không cần tune $\lambda$ quá kỹ, 0.5 là lựa chọn an toàn.

    Improvement không chỉ riêng DINOv2 — MAE và MoCov3 cũng cho
    thấy trend tương tự, chứng tỏ ý tưởng align representation
    là general, không phụ thuộc vào encoder cụ thể.
  }
\end{frame}

% ------------------------------------------------------------------
% APPENDIX SLIDE 8 – HYPERPARAMETERS & COMPUTE
% ------------------------------------------------------------------
\begin{frame}[t]
  \framesubtitle{Appendix}
  \frametitle{Hyperparameters and compute details}
  \begin{columns}[T,totalwidth=\textwidth]
    \column{0.60\textwidth}
    \begin{tabular}{L{2.8cm}L{3.8cm}}
      \toprule
      Category & Setup \\
      \midrule
      Backbone & Same DiT / SiT architecture \\
      Input & ImageNet 256\(\times\)256, VAE latents \(32{\times}32{\times}4\) \\
      Optimizer & AdamW, LR \(1{\times}10^{-4}\), no weight decay \\
      Batch size & 256 \\
      Sampler & Euler-Maruyama, \(w_t{=}\sigma_t\), NFE\(=\)250 \\
      Projection head & 3-layer MLP with SiLU \\
      REPA target & DINOv2-B/L (frozen) \\
      REPA depth & Block 8 (default for XL) \\
      \(\lambda\) & 0.5 \\
      Hardware & 8\(\times\) H100 80\,GB \\
      Throughput & \(\sim\)5.4 step/s \\
      \bottomrule
    \end{tabular}
  \column{0.34\textwidth}
    \callout{\textbf{Engineering note.} Training can be sped up further
    by \emph{precomputing} pretrained-encoder features offline, avoiding
    the DINOv2 forward pass at every step.}
    \vspace{0.18cm}
    \callout{\textbf{Overhead.} The REPA loss adds only the cost of a
    small MLP forward pass and one DINOv2 inference (or a lookup if
    precomputed). Negligible compared to the DiT/SiT forward--backward.}
  \end{columns}
  \note{
    Chi tiết hyperparameter. Input là ImageNet 256x256, encode bằng
    VAE thành latent 32x32x4. Batch size 256, optimizer AdamW
    với learning rate $10^{-4}$, không weight decay.

    Sampler: Euler-Maruyama với $w_t = \sigma_t$, 250 function
    evaluations khi sampling.

    Projection head: MLP 3 layer với SiLU activation. REPA target
    là DINOv2-B hoặc DINOv2-L (frozen). Align ở block 8.

    Hardware: 8 GPU H100 80GB, throughput khoảng 5.4 step/giây.

    Mẹo engineering: có thể precompute DINOv2 features offline
    cho toàn bộ dataset, tránh phải chạy DINOv2 forward mỗi step.
    Overhead của REPA rất nhỏ so với cost của DiT/SiT forward-backward.
  }
\end{frame}

% ------------------------------------------------------------------
% APPENDIX SLIDE 9 – DETAILED QUANTITATIVE RESULTS
% ------------------------------------------------------------------
\begin{frame}[t]
  \framesubtitle{Appendix}
  \frametitle{Detailed quantitative results}
  \begin{columns}[T,totalwidth=\textwidth]
    \column{0.48\textwidth}
    \textbf{Unconditional (no CFG)}
    \vspace{0.08cm}

    \begin{tabular}{L{2.4cm}C{0.8cm}C{0.9cm}C{0.9cm}}
      \toprule
      Model & Iter. & FID & Acc. \\
      \midrule
      SiT-B/2 + REPA & 400K & 24.4 & 61.2 \\
      SiT-L/2 + REPA & 400K & 10.0 & 69.4 \\
      SiT-XL/2 + REPA & 400K & 7.9 & 70.3 \\
      SiT-XL/2 + REPA & 4M & 5.9 & 74.6 \\
      \bottomrule
    \end{tabular}
    \vspace{0.16cm}
    {\small Larger models gain more from REPA, and the advantage
    keeps growing with longer training.}
  \column{0.48\textwidth}
    \textbf{With CFG and guidance interval}
    \vspace{0.08cm}

    \begin{tabular}{L{2.8cm}C{0.8cm}C{0.9cm}}
      \toprule
      Setting & \(w\) & FID \\
      \midrule
      Vanilla SiT-XL/2 & 1.50 & 2.06 \\
      REPA & 1.30 & 1.80 \\
      REPA + \([0, 0.75]\) & 2.00 & 1.78 \\
      REPA + \([0, 0.70]\) & 1.80 & \textbf{1.42} \\
      \bottomrule
    \end{tabular}
    \vspace{0.16cm}
    {\small Guidance interval plus REPA produces the best reported
    ImageNet 256\(\times\)256 result. The guidance weight \(w\) can be
    lower because the model already has strong semantics.}
  \end{columns}
  \note{
    Bảng kết quả chi tiết.

    Bên trái — không dùng CFG: SiT-B đạt FID 24.4 (nhỏ, ít cải thiện),
    SiT-L đạt 10.0, SiT-XL đạt 7.9 chỉ sau 400K iteration. Train thêm
    đến 4M iteration: 5.9.

    Điểm đáng chú ý: accuracy linear probe tăng theo model size và
    training time — chứng tỏ REPA cải thiện representation quality
    thực sự, không chỉ tối ưu cho FID.

    Bên phải — với CFG: vanilla SiT cần guidance weight $w = 1.5$ cho
    FID 2.06. REPA chỉ cần $w = 1.3$ cho FID 1.80 — guidance weight
    thấp hơn vì model đã có semantics tốt sẵn.

    Kết hợp guidance interval [0, 0.70] với $w = 1.8$: đạt 1.42 —
    state-of-the-art.
  }
\end{frame}

% ------------------------------------------------------------------
% APPENDIX SLIDE 10 – 512x512 EXTENSION
% ------------------------------------------------------------------
\begin{frame}[t]
  \framesubtitle{Appendix}
  \frametitle{Extension to ImageNet 512\(\times\)512}
  \begin{columns}[T,totalwidth=\textwidth]
    \column{0.60\textwidth}
    \AppFiveTwelve[width=\linewidth]
  \column{0.36\textwidth}
    \begin{itemize}
      \item Same REPA setup, larger latent input.
      \item REPA already beats vanilla SiT-XL/2 with
            \(\mathbf{>3\times}\) fewer iterations.
      \item Shows the method scales beyond the
            256\(\times\)256 headline experiment.
    \end{itemize}
    \vspace{0.14cm}
    \callout{The alignment idea is resolution-agnostic:
    it operates on patch tokens regardless of spatial size.}
  \end{columns}
  \note{
    Mở rộng sang 512x512. Cùng setup REPA, chỉ thay đổi input
    resolution. Latent lớn hơn (64x64x4 thay vì 32x32x4).

    Kết quả: REPA vẫn beat vanilla SiT-XL/2 với hơn 3 lần ít
    iteration hơn. Method scale tốt vì alignment hoạt động ở
    mức patch token — không phụ thuộc spatial size.

    Đây là bằng chứng rằng REPA không chỉ ``overfit'' cho
    benchmark 256x256.
  }
\end{frame}

% ------------------------------------------------------------------
% APPENDIX SLIDE 11 – TEXT-TO-IMAGE
% ------------------------------------------------------------------
\begin{frame}[t]
  \framesubtitle{Appendix}
  \frametitle{Text-to-image extension}
  \begin{columns}[T,totalwidth=\textwidth]
    \column{0.64\textwidth}
    \AppTTwoI[width=\linewidth]
  \column{0.32\textwidth}
    \begin{itemize}
      \item REPA applied to \textbf{MMDiT} on MS-COCO.
      \item ODE setup: FID improves from 6.05 to \textbf{4.73}.
      \item SDE setup: FID improves from 5.30 to \textbf{4.14}.
      \item Alignment remains useful even when
            \textbf{text conditioning} is present.
    \end{itemize}
    \vspace{0.14cm}
    \callout{REPA is not limited to class-conditional generation.
    The benefit transfers to richer conditioning modalities.}
  \end{columns}
  \note{
    Mở rộng sang text-to-image. Paper áp dụng REPA cho MMDiT
    (multi-modal diffusion transformer) trên MS-COCO.

    ODE sampling: FID giảm từ 6.05 xuống 4.73.
    SDE sampling: FID giảm từ 5.30 xuống 4.14.

    Đáng chú ý: khi đã có text conditioning (thông tin ngữ nghĩa
    phong phú từ text), REPA vẫn giúp ích. Tức là representation
    alignment bổ sung cho text conditioning, không phải redundant.

    Điều này mở ra khả năng áp dụng REPA cho các hệ thống T2I
    quy mô lớn hơn (Stable Diffusion, DALL-E scale).
  }
\end{frame}

% ------------------------------------------------------------------
% APPENDIX SLIDE 12 – FEATURE MAPS & QUALITATIVE
% ------------------------------------------------------------------
\begin{frame}[t]
  \framesubtitle{Appendix}
  \frametitle{Feature maps and extra qualitative samples}
  \begin{columns}[T,totalwidth=\textwidth]
    \column{0.62\textwidth}
    \AppFeature[width=\linewidth]
    \vspace{0.06cm}
    {\small Feature visualization shows REPA hidden states are more
    structured and semantically coherent, especially at early layers.}
  \column{0.34\textwidth}
    \AppQualTop[width=\linewidth]
    \vspace{0.14cm}
    \AppQualBottom[width=\linewidth]
    \vspace{0.06cm}
    {\small More generated samples showing consistent improvement
    across diverse ImageNet classes.}
  \end{columns}
  \note{
    Slide cuối appendix — visualization.

    Bên trái: feature map visualization. Hidden state của REPA
    có cấu trúc rõ ràng hơn, semantic coherent hơn — đặc biệt
    ở các layer sớm. Có thể thấy rõ biên vật thể, vùng nền
    được phân tách tốt hơn.

    Bên phải: thêm ảnh generate. Chất lượng nhất quán trên nhiều
    class khác nhau của ImageNet — cho thấy REPA không chỉ tốt
    cho một vài class cụ thể.

    Tổng kết: REPA là phương pháp đơn giản, hiệu quả, dễ tích
    hợp, và có potential mở rộng lớn. Cảm ơn mọi người!
  }
\end{frame}
